\section{Materiales Emisivos} % (fold)
\label{sec:materiales_emisivos}
La estructura de vóxeles utilizada durante el trazado de conos almacena un valor de radiancia por cada vóxel. Considerando esto, agregar materiales emisivos al proceso de voxelización es simple. Para la voxelización de materiales emisivos se utiliza otro volumen además de los ya existentes (albedo y normal) durante el proceso de voxelización. Este volumen almacena el promedio de emisión de los fragmentos que envuelve el vóxel. El valor de estos vóxeles es luego agregado al resultado de radiancia directa durante el proceso de sombreado de vóxeles. Los materiales emisivos pueden ser utilizados para aproximar cualquier clase de superficie luminosa como luces de área, luces de neón, magma, pantallas digitales, etc (ver Figura \ref{fig:emissive_mat}).
% section materiales_emisivos (end)

\begin{figure}[H]
	\centering
	\captionsetup{justification=centering}
	\includegraphics[width=0.95\linewidth]{media/emissive_neon.png}
	\caption{Materiales emisivos en Unreal Engine 4.6 \cite{unreal_46}.}
	\label{fig:emissive_mat}
\end{figure}